% ========================= ACT I: UNDERSTAND =========================
% storyboard.md Act I. Job: a sell, not a payoff -- version control is the
% prerequisite that unlocks Codex/Claude Code, not a demonstrated outcome.
% The mature, sophisticated result stays out of this act entirely (Act
% III's capstone).
%
% Task 3 is idea-first, not corpus-first: per Alex, "all of my proposal
% development and literature reviews are idea-first" -- a corpus enters
% after an AI-assisted refinement pass, not instead of starting from one.
% Grounded in three generic examples (chromatography, watershed,
% desalination).
%
% Slide 2 (LLM -> Chatbot -> Agent) deliberately does not repeat the
% Prologue's "Don't become loyal to a model" vendor table
% (01_ecosystem.tex) -- that slide is current offerings; this one is the
% conceptual distinction between kinds of tool.
%
% Task 1 (slide 3) does not mention emcal -- that example lives in Act
% II's Task 4 instead.

\stagedivider{2}{Act I: Understand}{Go from brainstorming or search with a chatbot to deeper engagement}

\begin{frame}{The filesystem is the source of truth}
  \vfill
  \begin{columns}[t]
    \begin{column}{0.46\textwidth}
      \begin{center}\textbf{A conversation}\end{center}
      \begin{itemize}
        \item not versioned
        \item not shared with collaborators
        \item not searchable next year
        \item gone when the session ends
      \end{itemize}
    \end{column}
    \begin{column}{0.46\textwidth}
      \begin{center}\textbf{A repository}\end{center}
      \begin{itemize}
        \item versioned
        \item shared by construction
        \item greppable
        \item outlives everyone involved
      \end{itemize}
    \end{column}
  \end{columns}
  \vfill
\end{frame}

\begin{frame}{An agent can touch your repository --- a chatbot can't}
  \vfill
  \begin{center}
  \begin{tikzpicture}[font=\small, node distance=1.1cm]
    \node[draw=NDBlue, rounded corners=2pt, minimum width=2.7cm, minimum height=1.3cm, align=center]
      (a) {\textbf{LLM}\\\scriptsize the model};
    \node[draw=NDBlue, rounded corners=2pt, minimum width=3.3cm, minimum height=1.3cm, align=center, right=of a]
      (b) {\textbf{Chatbot}\\\scriptsize LLM + a chat window\\\scriptsize text in, text out};
    \node[draw=NDGold, rounded corners=2pt, minimum width=3.6cm, minimum height=1.3cm, align=center, right=of b, fill=CodeBg]
      (c) {\textbf{Agent}\\\scriptsize LLM + tools\\\scriptsize reads and writes your files};
    \draw[-{Stealth[length=6pt]}, thick] (a) -- (b);
    \draw[-{Stealth[length=6pt]}, thick] (b) -- (c);
  \end{tikzpicture}
  \end{center}
  \vspace{0.6em}
  \begin{center}
    A chatbot only ever has the conversation. An agent can act on the repository you just built --- read the \texttt{CLAUDE.md}, edit the code, run the tests.
  \end{center}
  \vspace{0.4em}
  \begin{center}
    {\scriptsize The same tools from the roadmap slide, sorted by kind now, not by brand.}
  \end{center}
  \vfill
\end{frame}

\begin{frame}[fragile]{Task 1 --- set up git for a project}
  \vfill
  \begin{center}
  \begin{tikzpicture}[font=\small, node distance=1.2cm]
    \node[draw=OkVermillion, rounded corners=2pt, minimum width=3.4cm, minimum height=1cm, align=center]
      (a) {Local notebook\\\scriptsize no backup, no history};
    \node[draw=NDBlue, rounded corners=2pt, minimum width=2.6cm, minimum height=1cm, align=center, right=of a]
      (b) {\texttt{git init}};
    \node[draw=OkGreen, rounded corners=2pt, minimum width=3.4cm, minimum height=1cm, align=center, right=of b, fill=CodeBg]
      (c) {GitHub\\\scriptsize shared, backed up, versioned};
    \draw[-{Stealth[length=6pt]}, thick] (a) -- (b);
    \draw[-{Stealth[length=6pt]}, thick] (b) -- (c);
  \end{tikzpicture}
  \end{center}
  \vspace{0.5em}
  \begin{columns}[t]
    \begin{column}{0.31\textwidth}
      \centering\small\textbf{A commit}\\[0.2em]
      {\scriptsize a saved, timestamped snapshot, with a message explaining why}
    \end{column}
    \begin{column}{0.31\textwidth}
      \centering\small\textbf{A repository}\\[0.2em]
      {\scriptsize the full history --- any past state is one command away}
    \end{column}
    \begin{column}{0.31\textwidth}
      \centering\small\textbf{GitHub}\\[0.2em]
      {\scriptsize hosts it --- shared and backed up by construction}
    \end{column}
  \end{columns}
  \vspace{0.5em}
  \begin{center}
    \slidepoint{Version control is the \textbf{prerequisite} ---\\what it \textbf{unlocks} is the rest of this talk.}
  \end{center}
  \vspace{1.8cm plus 1fill}
\end{frame}

\begin{frame}{Task 2 --- a chat history can't hold any of these three}
  \vfill
  {\scriptsize A membrane-transport project's actual context, in each of the three kinds.}
  \vspace{0.4em}
  \begin{columns}[t]
    \begin{column}{0.31\textwidth}
      {\color{NDBlue}\textbf{Persistent}}\\[0.4em]
      Results are deterministic. A changed number is a finding, not noise.\\[0.7em]
      {\scriptsize$\rightarrow$ true across every session}
    \end{column}
    \begin{column}{0.31\textwidth}
      {\color{NDBlue}\textbf{Authoritative}}\\[0.4em]
      The actual model code and experimental data --- not a description of them.\\[0.7em]
      {\scriptsize$\rightarrow$ point at it, never summarize it}
    \end{column}
    \begin{column}{0.31\textwidth}
      {\color{NDBlue}\textbf{Task}}\\[0.4em]
      This run's question: replicate weighting, right now.\\[0.7em]
      {\scriptsize$\rightarrow$ the prompt (save it if it matters)}
    \end{column}
  \end{columns}
  \vfill
\end{frame}

\begin{frame}[fragile]{Task 2 --- write down what you'd otherwise re-explain every session}
  \vfill
  \begin{columns}[c]
    \begin{column}{0.62\textwidth}
      \begin{lstlisting}[basicstyle=\ttfamily\tiny]
# Project

Membrane transport parameter estimation.

## Where things are
  src/       the model
  results/   written by scripts; never by hand
  configs/   every run is a config file

## How to run
  conda activate diafiltration
  python -m pytest -m "not slow"

## Conventions
  Results are deterministic. A change in a
  number is a finding, not noise.

## Current state
  Working on: replicate weighting
  Resume at: notes/log.md, 2026-08-14
      \end{lstlisting}
    \end{column}
    \begin{column}{0.34\textwidth}
      Name it neutrally.\\[0.6em]
      \texttt{CLAUDE.md} and \texttt{AGENTS.md}\\
      become three-line pointers to it.\\[1.2em]
      Neither tool reads the other's file --- so the content lives
      in one place and the tool files just point.
    \end{column}
  \end{columns}
  \vspace{1.8cm plus 1fill}
\end{frame}

\begin{frame}{Task 3 --- idea first: brainstorm, refine, then bring in the literature}
  \vfill
  \begin{center}
  \begin{tikzpicture}[font=\scriptsize, node distance=0.4cm]
    \node[draw=NDBlue, rounded corners=2pt, minimum width=2.5cm, minimum height=1.3cm, align=center]
      (a) {Dictate the\\loose idea};
    \node[draw=NDBlue, rounded corners=2pt, minimum width=2.5cm, minimum height=1.3cm, align=center, right=of a]
      (b) {AI refines it\\report or prompt};
    \node[draw=NDBlue, rounded corners=2pt, minimum width=2.5cm, minimum height=1.3cm, align=center, right=of b]
      (c) {Corpus-grounded\\literature review};
    \node[draw=NDGold, rounded corners=2pt, minimum width=2.5cm, minimum height=1.3cm, align=center, right=of c, fill=CodeBg]
      (d) {Report or\\proposal};
    \draw[-{Stealth[length=5pt]}, thick] (a) -- (b);
    \draw[-{Stealth[length=5pt]}, thick] (b) -- (c);
    \draw[-{Stealth[length=5pt]}, thick] (c) -- (d);
  \end{tikzpicture}
  \end{center}
  \vspace{0.6em}
  \begin{center}
    Chromatography, watershed, desalination --- all three this same pipeline, generic. The corpus arrives to pressure-test an idea already forming, not to generate one.
  \end{center}
  \vspace{0.4em}
  \begin{center}
    {\scriptsize Tip: interpret the evidence, don't copy the abstract --- say what it changes about the next decision, not what each paper already says about itself.}
  \end{center}
  \vspace{1.8cm plus 1fill}
\end{frame}
